\section{The coupling as an energetic ruler}\label{sec:ruler}
The ruler is classical. Its five symbols are $W,V,A,S,D$: a horizontal
wall, a vertical divider, a site above the signal, the signal, and a
site below it. Allowed horizontal pairs, in left-to-right order, are
\[
 WW,\ AA,\ SA,\ DS,\ DD,\ AV,\ SV,\ VS,\ VD.
\]
Allowed vertical pairs, in top-to-bottom order, are
\[
 WS,\ WA,\ WV,\ SW,\ DW,\ VW,\ AA,\ AS,\ SD,\ DD,\ VV.
\]
Every other pair costs one. Each $W$ site receives energy $-g$;
each vertical $VW$ bond costs $r=1/16$. These are fixed rational
onsite and nearest-neighbour terms, affine in $g$.

\begin{lemma}[Forced cell geometry]\label{lem:forcing}
On a torus, every violation-free configuration is either all $A$, all
$D$, or a stack of stripes separated by $W$ rows. A stripe of height
$h\geq1$ consists of repeated columns
\[
 V,\ S_0,\ldots,S_{h-1},\qquad S_o=A^o S D^{h-1-o}.
\]
Its cells have pitch $p=h+1$, interior $h\times h$, and $p$ divides
the horizontal circumference.
\end{lemma}
\begin{proof}
The only horizontal pair containing $W$ is $WW$, so a wall occupies an
entire row. Vertical $WW$ is forbidden. Between walls a column
containing $V$ is all $V$. Any other column is $A^*SD^*$, $A^*$ or
$D^*$ by the vertical rules; the boundary rules forbid the last two.
If the signal in one column is at offset $o<h-1$, the pairs $SA$
and $DS$ force the next signal to offset $o+1$. At $o=h-1$ the next
column must be $V$: otherwise its last entry would be $A$ above $W$.
After $V$, the next signal has offset zero. This gives the cycle.
With no $W$, a cyclic vertical column is all $A$, all $D$ or all $V$.
The horizontal rules admit no cyclic mixture of these column types;
only the two uniform vacua remain.
\end{proof}

A perfect pitch-$p$ configuration has energy per site
\begin{equation}\label{eq:ruler-energy}
 f_g(p)=-\frac gp+\frac r{p^2},\qquad
 g_p=\frac{r(2p+1)}{p(p+1)},\qquad
 f_g(p)-f_g(p+1)=\frac{g_p-g}{p(p+1)}.
\end{equation}
Only $p\geq2$ is geometrically needed. For sufficiently small $g>0$
the unique minimizer is $p_*(g)=p$ on $(g_p,g_{p-1})$; at $g_p$ both
$p,p+1$ minimize. Completing the square in $1/p$ gives
\begin{equation}\label{eq:ruler-asymptotic}
 p_*(g)=\frac{2r}{g}+O(1),\qquad
 \min_p f_g(p)=-\frac{g^2}{4r}+O(g^4).
\end{equation}
The following estimate controls arbitrary defects, not only the perfect
patterns used in this minimization.

\begin{lemma}[Charging estimate]\label{lem:charging}
Put $f_* =\min(0,\min_{p\geq1} f_g(p))$. On an $L_1\times L_2$
torus let $V(\sigma)$ count violated pairs. For every column containing
$W$, let $s$ range over segments from one $W$ to the next, of pitch
$p_s$. Then
\begin{equation}\label{eq:charging}
 E(\sigma)-L_1L_2f_*
 \geq\sum_s p_s(f_g(p_s)-f_*)+(1-6r)V(\sigma).
\end{equation}
In particular a configuration with a violated pair lies at least $5/8$
above $L_1L_2f_*$.
\end{lemma}
\begin{proof}
A column with no $W$ has no negative bonus and contributes at least
$L_2f_*$. For the other columns assign each $W$ to the segment it starts.
If $v_s$ is the indicator that the segment ends in a $VW$ bond, its
energy contribution is $-g+rv_s$. The identity
\[
 -g-pf_*=p(f_g(p)-f_*)-r/p
\]
leaves the additional nonnegative residual $-A_0f_*$, where $A_0$
is the area of the $W$-free columns, and reduces the assertion to
$\sum_s(v_s-1/p_s)\geq-6V(\sigma)$.
Declare a segment vertically bad if its closed row range contains a
forbidden vertical bond. Group the remaining segments into maximal
horizontal runs having the same top and bottom wall rows and no
violation on the intervening horizontal bonds. Lemma~\ref{lem:forcing}
forces the column cycle inside a run. A run of $m$ columns and pitch
$p$ contains at least $\lfloor m/p\rfloor$ divider columns; its deficit
is greater than $-1$. A run closing around the torus has zero deficit.
Charge only the right end of each nonclosing run. A missing matching
wall segment produces a forbidden $W$--non-$W$ horizontal pair.
A matching but vertically bad segment supplies a forbidden vertical
bond; otherwise the failed join supplies a forbidden horizontal bond.
A violated pair belongs to the closed vertical ranges of at most three
segments on that side, and so pays for at most three such run ends.
A segment excluded from every good run contains a violated vertical
pair; its individual deficit is at least $-1$. Each such violation is
in at most three closed segment ranges. Charging both types uses at
most six units of deficit per violation. Substitution proves
\eqref{eq:charging}.
\end{proof}

\begin{corollary}[Selection and separation]\label{cor:selection}
For $g\leq0$ the vacua have energy zero and every other wall
configuration has energy at least $r$. For small positive nontie $g$,
on a square torus divisible by $p_*(g)$ the minimizers are exactly the
perfect pitch-$p_*$ cell lattices. Every other configuration lies at
least $\min(5/8,L\delta(g))$ above, where
\[
 \delta(g)=\min_{p\ne p_*}p(f_g(p)-f_*)>0.
\]
At ties allow both minimizing pitches and omit both from this minimum.
\end{corollary}
\begin{proof}
For $g\leq0$ a clean nonvacuum stack has energy
$L_1\sum_i(-g+r/p_i)\geq r$, since $p_i\leq L_1$.
Defective configurations obey Lemma~\ref{lem:charging}. For $g>0$,
clean stripes of a wrong pitch contribute $Lp(f_g(p)-f_*)$; all
other configurations have a violated pair. The minimum defining
$\delta$ is effective and positive because its terms tend to infinity
with $p$. The same argument excludes all pitches except the two
minimizers at a tie. Uniform vacua have the positive excess
$-L^2f_*$ and are also above the stated bound for sufficiently large $L$.
\end{proof}

The threshold for the last statement can be computed at rational
$g>0$: compare the rational energies, find the finite minimizing
pitches and a positive separation, then take $L=n!$ large enough that
these pitches divide $L$ and $L\delta\geq5/8$ (and the vacuum excess
also exceeds $5/8$). No assertion is made for incommensurate tori.
